\section{Introduction et contexte du stage}

Ce travail s’inscrit dans la continuité du projet eSWARM, dont l’objectif est de développer un système de drones quadricoptères modulaire et reconfigurable. Les appareils doivent être capables de s’assembler en vol afin d’accomplir des missions qu’un drone isolé ne pourrait réaliser, comme le transport de charges lourdes \cite{bonnel_modeling_2024}. Les systèmes robotiques modulaires et reconfigurables ont été définis dans \cite{seo_modular_2019, brooks_overview_2021}. Ils reposent sur une architecture composée de modules pouvant s’assembler dans différentes configurations et se réorganiser dynamiquement en fonction de la tâche à accomplir.

Parmi les projets de référence, le système ModQuad \cite{saldana_modquad_2018} illustre la faisabilité de telles approches en permettant à des quadricoptères de former des structures plus complexes par assemblage. La difficulté de ce type de système est la conception des lois de commande permettant de piloter une structure dont la configuration peut changer au cours du temps.

Dans ce contexte, et afin de réduire les contraintes expérimentales liées aux drones (autonomie limitée, risque de crash, fragilité du matériel), une alternative reposant sur des robots terrestres holonomes (pouvant se déplacer sans contrainte dans le plan) a été mise en place pour tester les algorithmes. L’objectif est ici de développer, implémenter et comparer différents algorithmes afin d'optimiser la commande d'un essaim de robots.
